% Parte do material do minicurso "Escrita de Artigos Científicos com LaTeX" (Conecta Hub)
% Exemplo 02 – Seções, listas e formatação de texto

\documentclass[12pt, a4paper]{article}
\usepackage[utf8]{inputenc}
\usepackage[T1]{fontenc}
\usepackage[brazil]{babel}

\title{Seções e Formatação}
\author{Seu Nome}
\date{}

\begin{document}

\maketitle

\section{Introdução}
Seções são criadas com \texttt{\textbackslash section\{Título\}}.

\subsection{Formatação de texto}
Use \textbf{negrito}, \textit{itálico}, \underline{sublinhado} e \texttt{código monoespaçado}.

Caracteres especiais precisam de escape: \% \& \$ \# \_ .

Aspas em português: ``texto entre aspas'' (use duas crases e duas aspas simples).

\subsection{Listas}

Lista com marcadores:
\begin{itemize}
    \item Primeiro item;
    \item Segundo item;
    \begin{itemize}
        \item Sub-item, com recuo.
    \end{itemize}
\end{itemize}

Lista numerada:
\begin{enumerate}
    \item Passo um;
    \item Passo dois;
    \item Passo três.
\end{enumerate}

Lista de descrição:
\begin{description}
    \item[LaTeX] sistema de composição de textos;
    \item[BibTeX] gerenciador de referências.
\end{description}

\section{Espaçamento e quebras}
Use uma linha em branco para novo parágrafo.\\
A barra dupla \texttt{\textbackslash\textbackslash} força uma quebra de linha.

\section{Conclusão}
Pronto: com seções e listas você já estrutura qualquer texto.

\end{document}
